\par\Needspace{4\baselineskip}
\begin{center}
\SourceLeaf{022b}\pdfbookmark[2]{4.}{article-004}4.
\end{center}
Anstelle des Abstands $x$ führen wir nun den Winkel $v$ zwischen der Apsidenlinie und der von der Sonne zum Ort des Körpers gezogenen Geraden (dem \emph{Radiusvektor}) ein. Dieser Winkel werde von dem Teil der Apsidenlinie aus gemessen, für den die Abstände $x$ positiv sind, und wachse in Richtung der Bewegung des Körpers. So ergibt sich $x=r\cos v$, und unsere Formel lautet damit $r=\frac{p}{1+e\cos v}$. Daraus können wir unmittelbar Folgendes schließen.

I. Für $v=0$ nimmt der Radiusvektor $r$ seinen kleinsten Wert an, nämlich $=\frac{p}{1+e}$. Dieser Punkt heißt \emph{Perihel}.

II. Entgegengesetzten Werten von $v$ entsprechen gleiche Werte von $r$. Daher teilt die Apsidenlinie den Kegelschnitt in zwei gleiche Teile.

III. In der \emph{Ellipse} wächst $r$ von $v=0$ an stetig, bis es den größten Wert $\frac{p}{1-e}$ im \emph{Aphel} bei $v=180^{0}$ erreicht. Nach dem Aphel nimmt es wieder so ab, wie es zuvor zugenommen hat, bis bei $v=360^{0}$ erneut das Perihel erreicht ist. Der Teil der Apsidenlinie, der auf der einen Seite vom Perihel und auf der anderen vom Aphel begrenzt wird, heißt \emph{große Achse}. Daher ist die große Halbachse, die auch \emph{mittlerer Abstand} heißt, $=\frac{p}{1-ee}$. Der Abstand des Mittelpunkts der Achse (des \emph{Mittelpunkts der Ellipse}) vom Brennpunkt ist $\frac{ep}{1-ee}=ea$, wobei $a$ die große Halbachse bezeichnet.

IV. Bei der \emph{Parabel} gibt es dagegen streng genommen kein Aphel. Der Wert von $r$ wächst unbegrenzt, je näher $v$ an $+180^{0}$ oder $-180^{0}$ herankommt. Für $v=\pm180^{0}$ wird der Wert von $r$ unendlich. Das zeigt, dass die Apsidenlinie die Kurve auf der dem Perihel gegenüberliegenden Seite nicht schneidet. Streng genommen kann man daher weder von einer großen Achse noch von einem Mittelpunkt der Kurve sprechen. Nach dem üblichen Vorgehen in der Analysis erweitern wir jedoch die für die Ellipse gefundenen Formeln, weisen der großen Achse einen unendlichen Wert zu und verlegen den Mittelpunkt der Kurve in unendliche Entfernung vom Brennpunkt.

\SourceLeaf{023}V. Bei der \emph{Hyperbel} schließlich ist $v$ auf noch engere Grenzen beschränkt, nämlich auf den Bereich zwischen $v = -(180^0-\psi)$ und $v = +(180^0-\psi)$. Dabei bezeichnet $\psi$ den Winkel mit dem Kosinus $= \frac{1}{e}\cdot$ Nähert sich nämlich $v$ diesen Grenzen, so wächst $r$ ins Unendliche. Setzte man für $v$ einen dieser Grenzwerte selbst ein, so ergäbe sich für $r$ ein unendlicher Wert. Das zeigt, dass eine Gerade, die mit der Apsidenlinie nach oben oder nach unten den Winkel $180^0-\psi$ bildet, die Hyperbel überhaupt nicht schneidet. Für die so ausgeschlossenen Werte, nämlich von $180^0-\psi$ bis $180^0+\psi$, weist unsere Formel $r$ einen negativen Wert zu. Die Gerade, die unter einem solchen Winkel zur Apsidenlinie geneigt ist, schneidet die Hyperbel selbst zwar nicht. Verlängert man sie aber nach rückwärts, so trifft sie den anderen Ast der Hyperbel. Dieser ist bekanntlich vom ersten Ast ganz getrennt und zu dem Brennpunkt hin konvex, in dem sich die Sonne befindet. Bei unserer Untersuchung, die, wie bereits bemerkt, auf der Annahme beruht, dass $r$ positiv anzusetzen ist, lassen wir diesen anderen Ast der Hyperbel jedoch außer Betracht. Auf ihm könnte sich nur ein Himmelskörper bewegen, auf den die Sonne nach denselben Gesetzen eine abstoßende statt einer anziehenden Kraft ausübte. — Streng genommen gibt es also auch bei der Hyperbel kein Aphel. Als Gegenstück zum Aphel kann man den Punkt des abgewandten Astes auffassen, der auf der Apsidenlinie liegt und den Werten $v = 180^0$, $r = -\frac{p}{e-1}$ entspricht. Wollen wir wie bei der Ellipse den Wert des Ausdrucks $\frac{p}{1-ee}$ auch hier, wo er negativ wird, als große Halbachse der Hyperbel bezeichnen, so gibt das Doppelte dieser Größe den Abstand des eben erwähnten Punkts vom Perihel an und zeigt zugleich, dass er auf der anderen Seite des Perihels liegt als der entsprechende Punkt bei der Ellipse. Ebenso nimmt $\frac{ep}{1-ee}$, also der Abstand des Mittelpunkts zwischen diesen beiden Punkten (des Mittelpunkts der Hyperbel) vom Brennpunkt, hier wegen seiner entgegengesetzten Lage einen negativen Wert an.
